\scanpage{141}\mspage{473}
tendue entre mes parents que tout allait être pour le mieux tant pour moi à \unsure{Blan|Blan-}\newline
kenese (près de Hambourg), que pour ma soeur qui depuis quelques mois avait été
larguée à la fin des fins dans une institution à Berlin pour enfants handicapés
(où on avait bien voulu d'elle, bien qu'elle n'était pas plus handicapée que moi
ou nos parents).

En dénouement à six mois étranges, lourds de menace sourde\quad et d'angoisse,
je me suis retrouvé du jour au lendemain dans un monde totalement différent du
seul monde que j'avais connu dans ma vie, celui formé par mes parents et par ma
soeur et moi. Je m'y retrouvais comme un parmi un groupe de pensionnaires, qui
mangions à part de la famille et faisions figure d'enfants de deuxième catégorie
pour les enfants de la maison, lesquels formaient un monde à part et nous regar-\newline
daient de haut. De ma mère je recevais une lettre hâtive et guindée de loin en
loin, et de mon père jamais une ligne de sa main, pendant les cinq ans que j'y
suis resté (jusqu'en 1939, à la veille de la guerre, quand j'ai fini par \unsure{rejoind|rejoindre}
mes parents sous la pression des événements).

Le couple qui m'avait accueilli m'a vite pris en affection. Aussi bien lui
ancien pasteur qui avait quitté le sacerdoce et vivait d'une maigre pension et
de leçons particulières de latin, de grec et de mathématiques, que sa femme pé-\newline
tillante de vie et parfois de malice, étaient des gens peu ordinaires, \unsure{attachant|attachants}
à bien des égards. Lui était un humaniste de vaste culture qui s'était un peu
égaré dans la politique, et avait eu maille à partir avec le régime nazi, qui a
fini par le laisser tranquille. Après la guerre j'ai renoué et je suis resté en
relations suivies avec eux jusqu'à la mort de l'un et de l'autre (*).

De lui et surtout d'elle, tout comme de mes parents, j'ai reçu du meilleur
comme aussi du pire. Aujourd'hui, avec un long recul, je leur suis reconnaissant
(comme je le suis à mes parents) pour ce ``meilleur'', comme aussi pour ce ``pire''.
C'est ce meilleur et ce pire que j'ai reçu, de mes parents d'abord, d'eux \unsure{ensuit|ensuite}
qui a formé le plus gros du volumineux ``paquet'' que j'ai reçu en partage dans \unsure{mo|mon}
enfance (comme chacun reçoit le sien...)\inserted{,}qu'il m'appartenait de déballer et \unsure{d'ex|d'exa-}\newline
miner. Ils font partie de la substance, de la richesse de mon passé, dont il ne
tient qu'à moi de nourrir mon présent.

Mon nouveau milieu était tout ce qu'il y a de ``comme il faut'' et de \unsure{confor|confor-}\newline
miste à beaucoup d'égards, avec en tous cas les attitudes répressives de rigueur
pour tout ce qui concerne le corps et, plus particulièrement, le sexe. Il a \unsure{fall|fallu}

\figure{\begin{tikzpicture}\draw(0,0)--(3,0);\end{tikzpicture}}
(*) Elle est morte à l'âge de 99 ans, il y a deux ans, et j'ai encore pu la voir
morte, en tête à tête avec elle, la veille de l'enterrement.
\mspage{140}
